\chapter{Summary and Conclusions}
\label{ch:conclusion}

The committee-requested proof of concept was achieved: a pICNN trained against
the compressible neo-Hookean reference supplied PyTorch-differentiated stress
and tangent to a nonlinear NGSolve solve through the integration-point bridge.
After those Phase~1 checks passed, the project extended the same workflow to a
periodic, mixed finite-strain RVE with exact rigid-particle kinematics.
Reaction-based homogenization and symbolic static condensation provided
work-conjugate macrostress and a full \(9\times9\) consistent tangent without
finite perturbation. Free particle rotation and the periodic translation gauge
were treated as distinct physical and null coordinates.

The verification programme produced near-machine-precision evidence for
periodicity, rigid motion, rotational equilibrium, Hill--Mandel consistency,
and tangent symmetry. Schema-2 production datasets preserve microscopic
Jacobian diagnostics and supported grouped training and independent numerical
evaluation. The Phase~6/7 checkpoint completed first-order macroscopic paths to
5\% stretch at three mesh sizes; the selected shear-adapted
checkpoint completed a matched first-order run at one mesh. These results
demonstrate substantial programming contributions in constitutive
modelling, finite elements, homogenization, data generation, and testing.

The final shear-adapted checkpoint was evaluated on a new direct-RVE set with
unseen shear amplitudes and a different microstructure seed. It reduced mean
relative shear stress and tangent errors by about 32\% and the full-stress
relative RMSE by about 28\% compared with the frozen Phase~6/7 model. A matched
first-order macro-FE calculation reached 5\% stretch with positive sampled
Jacobians. The remaining errors, approximately 40\% in the mean shear
stress/tangent measures, are still substantial. The independent set contains
only 12 scored cases on one new microstructure seed, so the generalization
claim is correspondingly limited.

Second-order mesh refinement is available only at 0.5\% stretch; a second-order
run to 5\% and wider quadrature convergence were not completed. The first-order
5\% energy comparison is matched on one mesh and does not establish response
convergence. The project therefore demonstrates a verified RVE-to-surrogate
workflow and a working macro-FE integration, but not a uniformly accurate
constitutive replacement or experimental validation. Further work should
expand independent microstructure and loading-path coverage, improve shear
accuracy, complete high-order load paths, and compare with experiments. No
experimental validation is claimed.
